\section{Design axes}
\label{sec:design}

Adding an entropic GW penalty is a design, not just an equation.
We identified five orthogonal choices --- one per Slurm job script
under \texttt{jobs/} --- and fix defaults that are informed by the
MMAudio base recipe. The defaults (leftmost entry in each axis) are
set in \texttt{config/base\_config.yaml} under the
\texttt{gw\_regularization} block. Every job overrides exactly the
axis it sweeps.

\subsection*{D1. Representation pair (variant)}
\emph{Which $v_i, a_j$ do we build $\DV, \DA$ from?}
The four variants are implemented in
\texttt{\_extract\_representations} in
\texttt{mmaudio/model/gw\_regularization.py}:
\begin{description}[leftmargin=2em]
  \item[\texttt{global}] (default). Raw CLIP features
    mean-pooled over time, $v = \overline{z^{\text{clip}}}$, vs
    raw $\ell_2$-normalized audio latent pooled over time,
    $a = \overline{x_1}$. No learned projection on either side.
  \item[\texttt{projected}] Apply MMAudio's own input projections
    --- \texttt{clip\_input\_proj} and \texttt{audio\_input\_proj} ---
    and mean-pool. Puts both modalities in the network's
    $D{=}448$ hidden space.
  \item[\texttt{c\_g}] Go one layer deeper on the video side: pool
    then apply \texttt{clip\_cond\_proj} to obtain the side that
    actually modulates the adaLN path. Compares the
    \emph{conditioning} vector that drives generation against the
    pooled audio latent. The audio side remains the projected
    pooling of $x_1$.
  \item[\texttt{fused}] Same representations as
    \texttt{projected}, but GW is replaced by FGW (cosine
    cross-cost $1 - \hat v \hat a^{\top}$, $\alpha=0.5$).
\end{description}

These variants differ along the \emph{how-deep-do-we-reach}
dimension: from raw pretrained features (\texttt{global}) down to
the specific vector that conditions generation (\texttt{c\_g}).
Constraining features deeper in the stack is more aggressive --- it
directly shapes generation --- but also more entangled with the
flow-matching objective, so harder to interpret.

\subsection*{D2. Penalty weight $\lambda_{\mathrm{GW}}$}
\emph{How strong is the relational pull?}
We sweep $\lambda_{\mathrm{GW}} \in \{0.001, 0.005, 0.01, 0.05,
0.1\}$ at the best variant from D1 (\texttt{jobs/train\_lambda.job},
\texttt{BEST\_VARIANT=<var>}). Default is $0.01$. Because the
per-distance-matrix mean-normalization (Sec.\,\ref{sec:method:reg})
removes feature-magnitude drift, these values are approximately
comparable across variants, though they are not exchangeable with a
$\lambda$ on an un-normalized GW loss in other papers.

\subsection*{D3. Detach direction}
\emph{Does the video side move, or only the audio side?}
With \texttt{detach\_video=True} (default), we
\texttt{.detach()} the video representation before building
$\DV$. Gradients then flow only through the audio side, meaning the
regularizer shapes audio geometry to match a fixed video geometry
for the current batch. With \texttt{detach\_video=False}, both
sides adapt. For \texttt{variant=global} there are no learned
parameters on the video side, so \texttt{detach\_video} is a
no-op; for \texttt{projected}, \texttt{c\_g}, and \texttt{fused},
gradients can push the video-side projections toward the audio
geometry, which risks shortcutting: the network could satisfy GW
by collapsing video representations. We therefore treat
\texttt{detach\_video=True} as the safer default and sweep both
(\texttt{jobs/train\_detach.job}, 2-element array).

\subsection*{D4. Schedule}
\emph{When does GW start biting?}
The flow-matching objective is ill-defined for random weights, so
a strong GW pull from step 0 can dominate and drag the network
toward a degenerate audio geometry before it even learns to
reconstruct. We implement three schedules in
\texttt{gw\_regularization.lambda\_schedule}
(\texttt{jobs/train\_schedule.job}):
\begin{description}[leftmargin=2em]
  \item[\texttt{constant}] $\lambda_{\mathrm{GW}}(t) = \lambda_{\mathrm{GW}}$
    for $t \ge w$, $0$ otherwise. $w=\texttt{warmup\_steps}$,
    default $1000$ (base config) but $50{,}000$ for the schedule
    sweep to make the warmup regime visible.
  \item[\texttt{linear\_rampup}] linear 0 $\to$ $\lambda_{\mathrm{GW}}$
    over $t \in [0, w]$, then constant.
  \item[\texttt{cosine\_anneal}] after warmup, decay back to 0
    over the remaining iterations, $0.5(1 + \cos(\pi\,(t-w)/(T-w)))$.
\end{description}
\texttt{constant} is the base-config default. The schedule has no
effect before $t = w$ (for \texttt{constant}) or during the ramp
(for the others), and \texttt{cosine\_anneal} explicitly removes
GW pressure in late training so the flow-matching loss owns the
final polishing phase.

\subsection*{D5. GW\,$\times$\,Synchformer interaction}
\emph{Does GW substitute for temporal-sync conditioning, or
stack with it?}
Synchformer features enter MMAudio through its own
\texttt{sync\_input\_proj} branch and are summed into the adaLN
context. If GW captures relational structure well, we might see
smaller gains from GW when Synchformer is on (saturation) and
larger gains when Synchformer is off (substitution); or GW and
Synchformer might be orthogonal and stack additively. The
$2 \times 2$ cell in \texttt{jobs/train\_gwxsync.job} runs:
\texttt{(sync\,on, gw\,off)} (pure baseline),
\texttt{(sync\,off, gw\,off)},
\texttt{(sync\,on, gw\,on)},
\texttt{(sync\,off, gw\,on)}. Sync-off is implemented with
\texttt{disable\_sync=true}, which replaces Synchformer features
with the learned \texttt{empty\_sync\_feat} slot inside the
training loop.

\paragraph*{OOD generalization (complementary).}
Orthogonal to D1--D5 but sharing a job script
(\texttt{jobs/train\_ood.job}) is a held-out-class test.
\texttt{experiments/make\_ood\_split.py} samples 30 VGGSound
classes uniformly at random (fixed seed), removes them from the
training TSVs, and keeps only those classes in the eval TSVs. If
the relational regularizer captures class-level structure rather
than within-class shortcuts, the GW variant should close some of
the OOD gap against a baseline trained on the same held-out split.

\paragraph{Orthogonality, not a factorial grid.}
We deliberately do not run the full $4 \times 5 \times 2 \times 3
\times 4 = 480$ factorial. For each axis we sweep conditional on
the \emph{current-best} setting of the previously-tested axes ---
D1 variants at default $\lambda$, then D2 $\lambda$ sweep at the
winning variant, etc. --- in the order D1, D2, D3, D4, D5. This is
cheaper and sufficient if the axes interact weakly, which D5
directly tests. The \texttt{BEST\_VARIANT} env var in the later
job scripts threads that choice through.
